%% LyX 2.4.4 created this file.  For more info, see https://www.lyx.org/.
%% Do not edit unless you really know what you are doing.
\documentclass[english]{article}
\usepackage[T1]{fontenc}
\usepackage[latin9]{inputenc}
\usepackage{enumitem}
\usepackage{amsmath}
\usepackage{amsthm}
\usepackage{amssymb}
\usepackage{geometry}
\geometry{verbose,tmargin=1in,bmargin=1in,lmargin=1in,rmargin=1in}
\PassOptionsToPackage{normalem}{ulem}
\usepackage{ulem}

\makeatletter

%%%%%%%%%%%%%%%%%%%%%%%%%%%%%% LyX specific LaTeX commands.
%% Because html converters don't know tabularnewline
\providecommand{\tabularnewline}{\\}

%%%%%%%%%%%%%%%%%%%%%%%%%%%%%% Textclass specific LaTeX commands.
\numberwithin{equation}{section}
\numberwithin{figure}{section}
\newlength{\lyxlabelwidth}      % auxiliary length 

%%%%%%%%%%%%%%%%%%%%%%%%%%%%%% User specified LaTeX commands.
\usepackage{fancyhdr}
\usepackage{lastpage}
\pagestyle{fancy}
\usepackage{enumitem}
\usepackage{ifthen}
%sets math fonts to be more readable
\everymath=\expandafter{\the\everymath\displaystyle}
%\DeclareMathSizes{10}{10}{10}{10}
\usepackage{multicol}
\usepackage{advdate}
% Added by lyx2lyx
\setlength{\parskip}{\smallskipamount}
\setlength{\parindent}{0pt}

\makeatother

\usepackage{babel}
\begin{document}
\renewcommand{\author}{Dr.\,James Ashe}

\newcommand{\classNum}{335}

\newcommand{\classSec}{A}

\newcommand{\nameType}{Name:}

\newcommand{\examType}{Homework}

\renewcommand{\date}{Due date: \AdvanceDate[14] \today} %%\today

%%First page's header%%%%%%%%%%%%%%%%%%%%%%
\fancypagestyle{firststyle} %%header settings for first pagr
{
	\fancyhead[L]{MTH \classNum{}(\classSec{}) \author{}\\
	\textbf{\examType{}} \date{}}
	\fancyhead[C]{\textbf{\nameType}}
	\setlength{\headsep}{14pt}
}
%%%%%%%%%%%%%%%%%%%%%%%%%%%%%%%%%%%%%%%%%%

%%Next page's header%%%%%%%%%%%%%%%%%%%%%%%
%\setlist{nolistsep} %eliminates space between list items
\lhead{MTH \classNum{}(\classSec{})} 
\chead{\textbf{\nameType}} 
\cfoot{\thepage{}~of~\pageref{LastPage}} 
\setlength{\headsep}{5pt} 
%%%%%%%%%%%%%%%%%%%%%%%%%%%%%%%%%%%%%%%%%%%

%%Various settings; see the preamble for other settings%%%%%
%%\setenumerate[0]{leftmargin=12pt,itemindent=0pt,labelwidth=0pt}
\renewcommand{\labelenumi}{\textbf{\arabic{enumi}.}}
\renewcommand{\labelenumii}{\textbf{\alph{enumii}.}}
\thispagestyle{firststyle}
%%%%%%%%%%%%%%%%%%%%%%%%%%%%%%%%%%%%%%%%%%

\textbf{Homework Problems. }\emph{Solutions should be written/typed
in numerical order and clearly labeled. Each page should have two
problems (except the last page if there are an odd number of problems).
Unsolved problems should be included but left blank.}
\begin{enumerate}
\item Write in tabular form:

\begin{enumerate}
\item $A=\{x\mid x\mbox{ is a letter in "correct"}\}$\\
$A=\{\mbox{c,o,r,e,t}\}$
\item $B=\{x\in\mathbb{R}\mid x\mbox{ is positive and }x\mbox{ is negative}\}$\\
$B=\emptyset$ since no number is negative \uline{and }positive. $0$
is\uline{ }not\uline{ }negative or positive.
\item $C=\{x\in\mathbb{Z}\mid x-2\leq5$\}\\
$C=\{\dots,4,5,6,7\}$
\end{enumerate}
\item Write these sets in a set-builder form:

\begin{enumerate}
\item Let $A=\{3,4,5,6,7,\dots\}$\\
$A=\{x\mid x\in\mathbb{Z}$ and $x\geq3\}$ 
\item Let $B=\{3,5,7,9,\dots\}$\\
$B=\left\{ x\in\mathbb{Z}\mid x\geq3\mbox{ and \ensuremath{x}}\mbox{ is odd}\right\} $
\item Let $C$ be the US presidents elected after 1999. \\
$C=\{x\mid x\mbox{ is a US president elected after 1999}\}$
\end{enumerate}
\item Find all the subsets of $A=\{x,y,z\}$.\\
\begin{tabular}{ccc}
$\{x\}$ & $\left\{ x,y\right\} $ & $\left\{ x,y,z\right\} $\tabularnewline
$\{y\}$ & $\left\{ x,z\right\} $ & $\left\{ \emptyset\right\} $\tabularnewline
$\{z\}$ & $\left\{ y,z\right\} $ & \tabularnewline
\end{tabular}
\item Prove that every nonempty subset has a proper subset.

\begin{proof}
Let $A\neq\emptyset$. $\emptyset\subseteq A$ and $A\neq\emptyset$.
So then $\emptyset$ is a proper subset of $A$, i.e., $\emptyset\subset A$.
\end{proof}
\item Prove that if $A$ is a subset of $\emptyset$, then $A=\emptyset$.

\begin{proof}
$\emptyset\subseteq A$ since every element in $\emptyset$ is also
in $A$, and by assumption $A\subseteq\emptyset$ by the hypothesis.
So then $A=\emptyset$. 
\end{proof}
\item Prove that $A=\{2,3,4,5\}$ is not a subset of $B=\{x\mid x\mbox{ is even}\}$.

\begin{proof}
$5$ is not even since $5=2\cdot2+1$ ($2$ does not evenly divide
$5$). So then $5\notin B$ which implies that $A\subseteq B$ since
$5\in A$. 
\end{proof}
\item Prove that $A$ and $B$ are subsets of $A\cup B$.

\begin{proof}
Let $a\in A$. Since $a\in A$, $a$ is an element of $A$ or $B$.
Thus $a\in A\cup B$ which implies that $A\subseteq A\cup B$. 

Similarly, we can show that $B\subseteq A\cup B$.
\end{proof}
\item Prove that $A=A\cup A$ and $A=A\cap A$.

\begin{proof}
$\left(A=A\cup A\right)$. By exercise 7, $A\subseteq A\cup A$. So
we need only show that $A\cup A\subseteq A$. Let $a\in A\cup A$.
Then $a\in A$ or $a\in A$, implying that $A\cup A\subseteq A$.
Thus $A\cup A=A$

$\left(A=A\cap A\right)$. Let $a\in A$. Then $a\in A$ and $a\in A$.
Thus $A\subseteq A\cap A$. Let $a\in A\cap A$. Then $a\in A$ and
$a\in A$ by definition, and this implies that $A\cap A\subseteq A$. 

Thus $A=A\cap A$.
\end{proof}
\item Prove that $A\cap B$ is a subset of $A$ and of $B$.

\begin{proof}
Let $x\in A\cap B$. Then $x\in A$ and $x\in B$ by definition. So
then $x\in A$ implying that $A\subseteq A\cap B$. Similarly we can
show that $B\subseteq A\cap B$.
\end{proof}
\item Prove that $A\cup B=\emptyset$ implies that $A=\emptyset$ and $B=\emptyset$.

\begin{proof}
Suppose that $A\cup B=\emptyset$. $A\subseteq A\cup B=\emptyset$
by exercise 7. Thus $A\subseteq\emptyset$, and this implies that
$A=\emptyset$ by exercise 5.

Similarly we can show that $B=\emptyset$.
\end{proof}
\end{enumerate}
\textbf{Presentation problems.}
\begin{enumerate}
\item Prove that $A-B$ is a subset of $A\cup B$.
\item Prove that if $A$ a subset of $B$ then $A\cap B=A$.
\item Let $A\cap B=\emptyset$. Prove that $B\cap A'=B$.
\item Prove that if $A$ is a subset of $B$ then $A\cup B=B$.
\item Prove that $A'-B'=B-A$.
\item Prove that if $A$ is a subset of $B$ then $B'$ is a subset of $A'$.
\item Prove that if $A\cap B=\emptyset$ then $A\cup B'=B'$.
\item Prove that $\left(A\cap B\right)'=A'\cup B'$.
\item Prove that $A\subset B$ implies that $A\cup\left(B-A\right)=B$.
\end{enumerate}
\vfill
\end{document}
